\section{Implementation}
\label{sec:implementation}

The design of the voltage comparator responsible for generating the digital
output code of the analog-to-digital converter was not arrived at directly.
It evolved through a sequence of candidate topologies, each evaluated against
the constraints imposed by the target supply voltage, the required conversion
accuracy, the effective input dynamic range, and the overall demands on circuit
compactness. The following subsections document this progression, presenting
each topology that was considered and explaining the physical reasoning that
either led to its adoption or prompted its replacement by a subsequent
candidate.


\subsection{First Candidate: Current-Mirror Flash Analog-to-Digital Converter}

The first topology investigated for the comparator core was the current-mirror
flash ADC (CMF-ADC) introduced by Lee~\textit{et al.}~\cite{lee2017},
originally proposed as the sensing front end of a coarse-fine dual-loop digital
low-dropout regulator. Although the application context of that work differs
substantially from the present design, the operating principle of the CMF-ADC
constitutes a direct voltage-to-thermometer-code conversion mechanism that is
directly applicable in the context of multi-bit comparison, making it a natural
starting point for the architectural exploration.

The circuit is divided into two cascaded functional stages: a transconductance
cell and a current-to-code converter. The transconductance cell is implemented
as a differential pair biased by a tail current of $2I_B$. Its function is to
translate the voltage difference between the input voltage, $V_\mathrm{in}$,
and an adjustable reference voltage, $V_\mathrm{ref,c}$, into an error
current $I_\mathrm{err}$, yielding branch currents
$I_1 = I_B - I_\mathrm{err}$ and $I_2 = I_B + I_\mathrm{err}$, where a
positive $I_\mathrm{err}$ corresponds to $V_\mathrm{in}$ exceeding
$V_\mathrm{ref,c}$. A third current, $I_3$, replicates $I_1$ through a
current mirror and is subsequently distributed to the conversion stage.

The current-to-code converter generates a multi-bit thermometer output by
replicating $I_2$ and $I_3$ at each comparison node through PMOS and NMOS
current mirrors, respectively, sized with deliberately asymmetric
multiplication factors. In the five-bit implementation described
in~\cite{lee2017}, the PMOS mirror ratios decrease monotonically from
$\times 14$ to $\times 6$ across the five output nodes specifically
$\times 14$, $\times 12$, $\times 10$, $\times 8$, and $\times 6$ for
bits $C_\mathrm{LPT}[4]$ through $C_\mathrm{LPT}[0]$ while the NMOS
mirror ratios increase correspondingly from $\times 6$ to $\times 14$.
At each comparison node, the output logic level is determined by whichever
mirrored current dominates: when the PMOS contribution exceeds the NMOS
contribution, the node is driven to a logic high; otherwise, the NMOS
pull-down prevails and the output settles to a logic low.

To explain the operating principle clearly, it is most instructive to consider
a configuration with an odd number of output bits, of which the five-bit case
is the natural choice. An odd bit count provides a well-defined, symmetric
midpoint whose comparison threshold aligns exactly with the condition
$V_\mathrm{in} = V_\mathrm{ref,c}$, making the midpoint of the thermometer
code unambiguous and facilitating a direct account of the circuit behaviour at
the reference voltage. When $V_\mathrm{in}$ lies substantially below
$V_\mathrm{ref,c}$, the error current $I_\mathrm{err}$ is strongly negative,
meaning $I_2 < I_B$ while $I_3 > I_B$. Since the NMOS mirrors replicate
$I_3$ with larger multiplication factors than the PMOS mirrors replicate $I_2$
at every node, the NMOS current dominates throughout and the output code
resolves to $00000$.

As $V_\mathrm{in}$ rises toward $V_\mathrm{ref,c}$, the magnitude of
$I_\mathrm{err}$ decreases and the comparison nodes with the highest
PMOS-to-NMOS mirror ratio are the first to transition. The node
$C_\mathrm{LPT}[4]$, carrying a PMOS mirror of ratio $\times 14$ against an
NMOS mirror of only $\times 6$, presents the most favourable balance in favour
of the PMOS side and therefore requires only a modest positive shift in $I_2$
relative to $I_3$ before its output transitions high. Node $C_\mathrm{LPT}[3]$
($\times 12$ PMOS, $\times 8$ NMOS) follows at a moderately larger error
current, while the least significant nodes, $C_\mathrm{LPT}[1]$ and
$C_\mathrm{LPT}[0]$, carry progressively more dominant NMOS mirrors and
require a correspondingly larger current excess before their outputs transition.

At the exact condition $V_\mathrm{in} = V_\mathrm{ref,c}$, the error current
vanishes and both $I_2$ and $I_3$ are equal to $I_B$. Nodes
$C_\mathrm{LPT}[4]$ and $C_\mathrm{LPT}[3]$ deliver PMOS currents of
$14 I_B$ and $12 I_B$, respectively, against NMOS contributions of $6 I_B$
and $8 I_B$, so both outputs are high. Nodes $C_\mathrm{LPT}[1]$ and
$C_\mathrm{LPT}[0]$ exhibit the reverse relationship and remain low. The
central node $C_\mathrm{LPT}[2]$ carries equal PMOS and NMOS currents of
$10 I_B$ and is therefore in an exactly balanced. The resulting thermometer code
at $V_\mathrm{in} = V_\mathrm{ref,c}$ is consequently $11{?}00$, where the
question mark denotes the indeterminate state of the central bit, fully
consistent with the expected behaviour of a comparator operating precisely at
its decision threshold.

The thresholds governing the transition of each bit are fully determined by
the mirror ratios, as established analytically in~\cite{lee2017}, and result
in a monotonic, symmetric thermometer-coded response across the input voltage
range. This architecture combines inherently fast conversion speed stemming
from the absence of a conventional error amplifier, whose finite bandwidth
would otherwise limit the response time of the control
loop~\cite{lee2017} with a straightforward, hardware-efficient structure and
an explicitly quantifiable input-output characteristic. It was therefore
adopted as the first candidate for the present comparator design, and its
operating principle informed all subsequent topological choices.


\subsection{Second Candidate: Complementary Dual-Input Architecture}

A notable limitation of the CMF-ADC topology, as configured, is that its
input common-mode range is restricted by the headroom requirements of the
differential pair and tail current source. In order to extend the effective
operating range of the comparator toward full rail-to-rail input coverage, a
complementary dual-input architecture was developed and evaluated as a second
candidate.

This topology employs two parallel transconductance stages operating
simultaneously: one based on an NMOS differential input pair and a second
based on a PMOS differential input pair. The motivation for this arrangement
is well established in the context of rail-to-rail input
amplifiers~\cite{razavi}: an NMOS-input stage operates correctly when the
common-mode input voltage is in the vicinity of the negative supply rail,
where the gate overdrive of the input transistors is sufficient to maintain
saturation, while a PMOS-input stage provides coverage near the positive
supply rail. Operated jointly, the two stages guarantee that at least one
transconductor remains fully functional across the complete supply voltage
span.

In the implementation explored here, the two transconductors were not combined
through a simple parallel superposition of their output currents. Rather, one
output branch of each OTA was directed toward the PMOS current mirrors
responsible for generating the individual bits of the thermometer code,
following the same current-comparison principle established in the CMF-ADC.
The remaining output branch of each OTA was cross-coupled to the complementary
transconductor: the residual output current from one stage was used to modulate
the output node of the other, and the same arrangement was applied
symmetrically in the reverse direction. This bidirectional cross-coupling was
intended to exploit the complementary transconductance to extend the effective
dynamic range of the conversion, allowing each stage to reinforce the other at
input voltages near the boundary of its individual valid operating region.

Despite the functional merit of this scheme, its practical realisation proved
incompatible with the target supply voltage of $V_\mathrm{DD} = 0.8$~V.
A standard five-transistor OTA requires a minimum headroom that accommodates
the gate-to-source voltage drop of the differential pair, the saturation
voltage of the tail current source, and the voltage drop across the active
load. At $0.8$~V, this budget is already severely constrained. The
cross-coupled dynamic range extension mechanism required one additional
transistor in the critical signal path beyond the five-transistor baseline,
and this single additional device was sufficient to push the cumulative
headroom demand beyond the available supply, preventing the circuit from
maintaining all active elements in saturation over the intended input range
and thus rendering the architecture non-functional in practice.

In an effort to recover the lost headroom margin, the conventional current
mirror structures within the OTA branches were replaced by the low-voltage
current mirror topology described by Razavi~\cite{razavi}. This configuration
reduces the minimum voltage required across the mirror branch by eliminating
one stacked gate-to-source drop from the critical headroom path, at the cost
of a moderate increase in circuit complexity. However, even with this
substitution, the cumulative voltage demand of the complete circuit 
encompassing the differential input pair, the tail current source, the
low-voltage mirror, and the cross-coupled output branch remained in excess
of $0.8$~V. No viable bias point could be identified that simultaneously
maintained all transistors in their correct operating region, and the
architecture was consequently abandoned.


\subsection{Adopted Architecture: Five-Transistor OTA with Feedback}

Following the rejection of the complementary dual-input approach, the design
converged on a classical five-transistor operational transconductance amplifier
as the foundation for the comparator. In this topology, the differential pair
is loaded by an active current mirror formed by two PMOS transistors: one
connected in diode configuration, with its gate tied directly to its own
drain, and a second transistor whose gate is driven by that same node. This
arrangement, in which the diode-connected device sets the mirror reference and
the second transistor replicates the resulting current to the complementary
output branch, constitutes the standard active load of a differential pair
and is the defining structural element of the five-transistor
OTA~\cite{razavi}. By dispensing with the additional stacked elements required
by the previous candidate, this topology fits within the available headroom at
$V_\mathrm{DD} = 0.8$~V and operates correctly across the intended input
range.

This topology also differs fundamentally from the current-mirror flash ADC in
the mechanism by which the comparison is realised. Rather than relying on a
separate current-to-code converter with asymmetrically ratioed mirror branches,
the five-transistor OTA performs a direct high-gain differential-to-single-ended
conversion, producing an output voltage that responds strongly to the sign and
magnitude of the differential input. To maximise the sensitivity of the
comparator to small input voltage differences, the OTA was designed to achieve
a high open-loop voltage gain.

% [TO BE COMPLETED: sizing or bias strategy -- W/L ratios, gm target, etc.]

This high gain, however, introduced a significant practical difficulty that
became apparent during circuit evaluation. Even a small unintended differential
voltage at the input arising from device offset, thermal noise, or a marginal
overdrive condition was sufficient to drive the output rapidly into hard
saturation before the circuit could respond to the intended comparison signal.
Once the output reached the supply rail or ground, recovery was limited by the
ability of the bias current to discharge or charge the high-impedance output
node, and this prolonged saturation degraded both the effective comparison
speed and the accuracy of successive conversions.

To resolve this behaviour, a feedback mechanism was incorporated around the
OTA.

% [TO BE COMPLETED: describe the specific feedback topology -- resistive,
% capacitive, regulated bias, or other -- including its effect on the
% closed-loop gain and the improvement in dynamic behaviour.]

The role of the feedback is to confine the output operating point to a
controlled linear regime, limiting the maximum output voltage excursion and
thereby preventing premature saturation, while preserving sufficient
differential gain for accurate comparison at the target conversion resolution.
